\begin{ouverture}[algorithmique]{21}{Instructions, répétitions et chemins}{Partie 4 · Données, hasard et instructions}
{Entrée · Sortie · Séquence · Répétition · Chemin}
\objectif{Identifier une instruction, une séquence, une entrée et une sortie.}
\objectif{Exécuter et produire une séquence d’instructions, y compris répétée.}
\objectif{Programmer un chemin simple, avec ou sans machine.}
\prerequis{Déplacements codés ; programmes de calcul ; motifs évolutifs.}
\end{ouverture}
\pagestyle{otMain}\essentiel\fontsize{11}{15}\selectfont
\phantomsection\label{cours-21-1}
\essSection{1}{Instruction, séquence, entrée et sortie}
\begin{definition}Une instruction est un ordre élémentaire donné à une personne ou à une machine : « avancer de 2 cases », « ajouter 3 ». Une séquence d’instructions est une suite d’instructions exécutées dans l’ordre. Une entrée est une donnée fournie au programme ; une sortie est un résultat qu’il affiche.\end{definition}
\label{fig-cours-21-1-cours-blocs}\figureGeo{cours-blocs}

\essSection{2}{Exécuter dans l’ordre et suivre les étapes}
\begin{savoirfaire}[Exécuter à la main]On note la valeur obtenue après chaque instruction. Avec l’entrée 5 : $5\times3=15$, puis $15+4=19$. La sortie est 19.\end{savoirfaire}
\begin{attention}L’ordre des instructions compte : « ajouter 4, puis multiplier par 3 » donne, avec l’entrée 5, $(5+4)\times3=27$, et non 19.\end{attention}
\clearpage
\phantomsection\label{cours-21-2}
\essSection{3}{Reconnaître et répéter une séquence}
\begin{definition}Une répétition, ou boucle, exécute plusieurs fois la même séquence. « Répéter 4 fois : avancer de 3 cases, tourner d’un quart de tour à gauche » trace un carré de 3 cases de côté.\end{definition}
\begin{savoirfaire}[Suivre une boucle]Pour répéter à la main, on note après chaque tour de boucle la position et l’orientation du personnage. On s’arrête après le nombre de tours annoncé.\end{savoirfaire}
\essSection{4}{Coder un chemin simple}
\label{fig-cours-21-2-cours-chemin}\figureGeo{cours-chemin}

\begin{savoirfaire}[Coder un chemin]On découpe le chemin en morceaux : avancer d’un certain nombre de cases, tourner d’un quart de tour à gauche ou à droite. Quand un même morceau revient plusieurs fois, on l’écrit avec une répétition.\end{savoirfaire}

\clearpage\rubrique{Découverte}
\begin{blocLivre}{ex-21-1}\begin{exercice}[1]{Instruction ou séquence ?}
Classer en instructions et en séquences d’instructions : « avancer de 3 cases » ; « avancer de 1 case, tourner à gauche, avancer de 2 cases » ; « afficher le résultat » ; « ajouter 5, puis multiplier par 2 ».
\end{exercice}\end{blocLivre}
\begin{blocLivre}{ex-21-2}\begin{exercice}[1]{Exécuter un programme}
Exécuter le programme avec l’entrée 7, en notant la valeur après chaque bloc. Quelle est la sortie ?
\label{fig-ex-21-2}\figureGeo{ins-2}
\end{exercice}\end{blocLivre}
\begin{blocLivre}{ex-21-3}\begin{exercice}[1]{Entrée et sortie}
Une machine presse des oranges : on lui donne le nombre d’oranges, chaque orange donne 8 cL de jus, puis la machine affiche la quantité de jus obtenue. Quelle est l’entrée ? La sortie ? Que s’affiche-t-il pour 6 oranges ?
\end{exercice}\end{blocLivre}
\begin{blocLivre}{ex-21-4}\begin{exercice}[1]{Suivre un chemin}
Le robot part de la case D, tourné vers la droite. Il exécute : avancer de 3 cases ; tourner à gauche ; avancer de 2 cases ; tourner à droite ; avancer de 2 cases. Sur quelle case, P, Q ou R, arrive-t-il ?
\label{fig-ex-21-4}\figureGeo{ins-4}
\end{exercice}\end{blocLivre}
\begin{blocLivre}{ex-21-5}\begin{exercice}[1]{L’ordre compte}
Avec l’entrée 5, exécuter le programme « ajouter 3, puis multiplier par 2 », puis le programme « multiplier par 2, puis ajouter 3 ». Comparer.
\end{exercice}\end{blocLivre}

\rubrique{Entraînement}
\begin{blocLivre}{ex-21-6}\begin{exercice}[2]{Répéter à la main}
Le lutin part du point de départ, tourné vers la droite. Exécuter « répéter 3 fois : avancer de 2 cases, tourner d’un quart de tour à droite ». Où arrive-t-il, et dans quelle direction regarde-t-il ?
\label{fig-ex-21-6}\figureGeo{ins-6}
\end{exercice}\end{blocLivre}
\begin{blocLivre}{ex-21-7}\begin{exercice}[2]{Écrire avec une boucle}
Réécrire avec une répétition la séquence : avancer de 2 cases, tourner à gauche, avancer de 2 cases, tourner à gauche, avancer de 2 cases, tourner à gauche, avancer de 2 cases, tourner à gauche. Quelle figure obtient-on ?
\end{exercice}\end{blocLivre}
\begin{blocLivre}{ex-21-8}\begin{exercice}[2]{Coder un chemin}
Écrire un programme qui conduit le robot de la case D à la case A sans passer par les cases grisées. Il part tourné vers la droite.
\label{fig-ex-21-8}\figureGeo{ins-8}
\end{exercice}\end{blocLivre}
\begin{blocLivre}{ex-21-9}\begin{exercice}[2]{Remonter un programme}
Un programme multiplie l’entrée par 4, puis ajoute 2. La sortie est 50. Quelle était l’entrée ?
\end{exercice}\end{blocLivre}
\begin{blocLivre}{ex-21-10}\begin{exercice}[2]{Le motif répété}
Le lutin part du point de départ, tourné vers la droite. Exécuter « répéter 4 fois : avancer de 1 case, tourner à gauche, avancer de 1 case, tourner à droite ». Quelle figure trace-t-il ? Où arrive-t-il ?
\label{fig-ex-21-10}\figureGeo{ins-10}
\end{exercice}\end{blocLivre}

\rubrique{Problèmes}
\begin{blocLivre}{ex-21-11}\begin{exercice}[3]{Le robot livreur}
Le robot part de D, tourné vers la droite, et exécute : avancer de 2 ; tourner à gauche ; avancer de 3 ; tourner à droite ; avancer de 4 ; tourner à droite ; avancer de 1. Devant quelle maison s’arrête-t-il ?
\label{fig-ex-21-11}\figureGeo{ins-11}
\end{exercice}\end{blocLivre}
\begin{blocLivre}{ex-21-12}\begin{exercice}[3]{Le polygone}
Un lutin exécute « répéter 6 fois : avancer de 3 cm, tourner de $60^\circ$ à gauche ». Quelle figure trace-t-il ? Quel est son périmètre ?
\label{fig-ex-21-12}\figureGeo{ins-12}
\end{exercice}\end{blocLivre}
\begin{blocLivre}{ex-21-13}\begin{exercice}[3]{Le tableur}
Dans un tableur, la cellule A1 contient 5 et la cellule A2 contient la formule « =A1+4 », recopiée vers le bas. Quel nombre s’affiche en A10 ? Dans quelle cellule lit-on 101 ?
\end{exercice}\end{blocLivre}
\begin{blocLivre}{ex-21-14}\begin{exercice}[3]{Trouver l’erreur}
Pour tracer un carré, Hugo a écrit : « répéter 3 fois : avancer de 4 cases, tourner à gauche ». Observer le tracé obtenu. Corriger le programme.
\label{fig-ex-21-14}\figureGeo{ins-14}
\end{exercice}\end{blocLivre}
\begin{blocLivre}{ex-21-15}\begin{exercice}[3]{Le plus court chemin}
Le robot doit aller de D à A, tourné vers la droite au départ, sans traverser les cases grisées. Il avance case par case. Combien de cases parcourt-il au minimum ? Écrire un programme court avec une répétition.
\label{fig-ex-21-15}\figureGeo{ins-15}
\end{exercice}\end{blocLivre}

\clearpage
